When AI agents share channels and credentials, one agent can press another to break its operator's rule, a first step towards collusion.
Output filters check only what an agent has already written.
Can an overseer tell beforehand that an agent is about to give in?
We show that a linear probe on the agent's activations can, in a controlled multi-agent setting.
Across six scenarios, a peer presses a target for up to 20 rounds to reveal a secret, such as a deploy token; the target often refuses, then gives in.
Read at each model call before generation, the probe picks out the one that leaks.
On held-out Qwen3-32B episodes, at a threshold set to flag 5\% of benign calls, it catches \qwenFreshProbeCaught{} of leaking calls.
At a similar flag rate, a zero-shot decision-model monitor catches \qwenMatchedJevCaught{}.
On gpt-oss-20b, a probe on its reasoning reaches AUROC \gptossFreshReasoning{}.
A bag-of-words classifier trained on the same calls comes close, so most of the gain comes from training on labels, not activations.
The probe also flags permitted disclosures: it reads the coming disclosure more than the violation.
The warning is short-range, and halting on any alarm would stop most benign episodes, which a second text check does not fix.
A warning raised before generation could let an overseer stop an agent before it gives in, before others act on its messages or recruitment spreads through a swarm.
